\documentclass[11pt]{article}

\usepackage[a4paper,margin=1in]{geometry}
\usepackage[T1]{fontenc}
\usepackage[utf8]{inputenc}
\usepackage{lmodern}
\usepackage{microtype}
\usepackage{amsmath,amssymb,amsthm,mathtools}
\usepackage{booktabs}
\usepackage{graphicx}
\usepackage{enumitem}
\usepackage{xcolor}
\usepackage[hidelinks]{hyperref}

\newtheorem{theorem}{Theorem}[section]
\newtheorem{proposition}[theorem]{Proposition}
\newtheorem{lemma}[theorem]{Lemma}
\newtheorem{corollary}[theorem]{Corollary}
\newtheorem{remark}[theorem]{Remark}
\theoremstyle{definition}
\newtheorem{definition}[theorem]{Definition}

\newcommand{\K}{\mathbb K}
\newcommand{\R}{\mathbb R}
\newcommand{\T}{\mathbb T}
\newcommand{\whotpi}{\widehat{\otimes}_{\pi}}
\newcommand{\co}{\operatorname{co}}
\newcommand{\supp}{\operatorname{supp}}
\newcommand{\Ree}{\operatorname{Re}}
\newcommand{\SDP}{\mathrm{SDP}}
\newcommand{\ccs}{\mathrm{ccs}}

\title{Super Daugavet Points in Projective Tensor Products\\
\large Verification, Literature Status, and Further Transfer Results}
\author{Automatic Math Discovery Run \texttt{fa\_banach\_001}}
\date{June 22, 2026}

\begin{document}
\maketitle

\begin{abstract}
The supplied argument proves a correct positive subcase of the question whether
super Daugavet points pass to projective tensor products.  Its key reduction is
that a one-dimensional $L$-summand in one tensor factor turns the projective
tensor product into an $\ell_1$-sum.  The underlying two-coordinate
$\ell_1$-stability fact is already published, so the argument is best viewed as
a clean tensor corollary rather than an independent resolution.

I did not locate a proof or a counterexample to the general question in the
literature available through 21 June 2026.  The report records one later
published enlargement and proves three further positive statements.  A 2026
characterization of super Daugavet points in vector-valued $L_1$ spaces already
implies that, for every $g\in S_{L_1(\mu)}$, the tensor $x\otimes g$ is super
Daugavet whenever $x$ is.  We give a self-contained arbitrary-$\ell_1$-sum
stability theorem, prove tensor transfer over both real and complex scalars for
the stronger class of \emph{ccs Daugavet points}, and prove the general
super-Daugavet tensor conclusion locally at every elementary target tensor.
The last result isolates the remaining obstruction: a weakly open subset of a
projective tensor ball may contain no elementary tensor at all.
\end{abstract}

\section{The question and the outcome}

Let $X$ be a Banach space and $x\in S_X$.  The point $x$ is a
\emph{super Daugavet point} if
\[
  \sup_{u\in W}\|x-u\|=2
\]
for every nonempty relatively weakly open subset $W$ of $B_X$.
Haller, Langemets, Perreau and Veeorg asked whether
\begin{equation}\label{eq:main-question}
  x\in S_X,\ y\in S_Y\text{ super Daugavet}
  \quad\Longrightarrow\quad
  x\otimes y\in S_{X\whotpi Y}\text{ super Daugavet}.
\end{equation}
They had already proved that a super $\Delta$-point tensors with an arbitrary
unit vector to give a super $\Delta$-point, while Dantas, Jung and Rueda Zoca
had proved the corresponding elementary-tensor transfer for ordinary
Daugavet points.

\begin{center}
\includegraphics[width=\textwidth]{figures/open_problem_crop.png}

\small Source crop from Haller--Langemets--Perreau--Veeorg, recording the
open super-Daugavet tensor-transfer sentence.
\end{center}

\medskip
\noindent
\textbf{Outcome of the present investigation.}
\begin{enumerate}[label=(\roman*),leftmargin=2.2em]
\item The supplied theorem and proof are correct.
\item The main ingredient---stability of super Daugavet points under
      $\ell_1$-sums---was already known in two-coordinate form.  Thus the
      supplied tensor statement is essentially an immediate application of
      known results and the standard tensor identity.
\item No full proof or counterexample to \eqref{eq:main-question} was found in
      the literature searched through 21 June 2026.
\item The positive results proved below substantially enlarge the supplied
      subcase, but they do not settle \eqref{eq:main-question} in full.
\end{enumerate}

All results below are valid over both the real and complex scalar fields.  In
the complex case, slices are defined using real parts of functionals.

\section{A net criterion}\label{sec:net}

We use repeatedly the following standard characterization.

\begin{lemma}[Net criterion]\label{lem:net}
Let $x\in S_X$.  The following are equivalent.
\begin{enumerate}[label=(\alph*)]
\item $x$ is a super Daugavet point.
\item For every $u\in B_X$ there is a net $(u_\alpha)\subset B_X$ such that
\[
  u_\alpha\xrightarrow{w}u
  \qquad\text{and}\qquad
  \|x-u_\alpha\|\longrightarrow 2.
\]
\end{enumerate}
\end{lemma}

\begin{proof}
Assume (a).  Direct the pairs $(U,\varepsilon)$, where $U$ is a relatively
weak neighborhood of $u$ in $B_X$ and $\varepsilon>0$, by refinement of $U$
and decrease of $\varepsilon$.  Choose
$u_{U,\varepsilon}\in U$ with
$\|x-u_{U,\varepsilon}\|>2-\varepsilon$.  This is the required net.

Conversely, let $W$ be a nonempty relatively weakly open subset of $B_X$ and
choose $u\in W$.  A net as in (b) is eventually in $W$, and therefore
$\sup_{v\in W}\|x-v\|=2$.
\end{proof}

The same proof gives the weak-star version in a dual space.

\section{Audit of the supplied partial result}\label{sec:audit}

The supplied theorem is the following.

\begin{theorem}[Supplied $L$-summand subcase]\label{thm:supplied}
Let $X$ and $Y$ be Banach spaces over $\K$.  Suppose that $x_0\in S_X$ is a
super Daugavet point and that
\[
       Y=\K y_0\oplus_1 Z
\]
isometrically for some $y_0\in S_Y$.  Then $x_0\otimes y_0$ is a super
Daugavet point of $X\whotpi Y$.
\end{theorem}

\begin{proof}
The projective tensor product distributes isometrically over $\ell_1$-sums:
\[
 X\whotpi(\K y_0\oplus_1 Z)
 \cong (X\whotpi\K y_0)\oplus_1(X\whotpi Z)
 \cong X\oplus_1(X\whotpi Z).
\]
Under this isometry, $x_0\otimes y_0$ corresponds to $(x_0,0)$.  It remains
to verify that $(x_0,0)$ is super Daugavet in $X\oplus_1 E$ for every Banach
space $E$.

Fix $(u,v)\in B_{X\oplus_1 E}$, put $s=\|v\|$ and $r=1-s$, and note that
$\|u\|\le r$.  If $r=0$, then $u=0$ and the constant net $(0,v)$ is at
distance $1+\|v\|=2$ from $(x_0,0)$.  If $r>0$, Lemma~\ref{lem:net} gives a
net $(p_\alpha)\subset B_X$ such that
\[
 p_\alpha\xrightarrow{w}u/r,
 \qquad \|x_0-p_\alpha\|\longrightarrow2.
\]
Set $w_\alpha=(rp_\alpha,v)$.  Then $w_\alpha\in B_{X\oplus_1 E}$,
$w_\alpha\to(u,v)$ weakly, and
\begin{align*}
 \|(x_0,0)-w_\alpha\|
 &=\|x_0-rp_\alpha\|+s\\
 &\ge \|x_0-p_\alpha\|-(1-r)\|p_\alpha\|+s\\
 &\ge \|x_0-p_\alpha\|.
\end{align*}
The left side is at most $2$, so it converges to $2$.  Lemma~\ref{lem:net}
finishes the proof.
\end{proof}

\subsection*{Audit verdict}
Every mathematical step above is valid.  In particular:
\begin{itemize}[leftmargin=1.8em]
\item the net characterization is correct;
\item the scaling $u/r$ is legitimate because $\|u\|\le r$;
\item the displayed distance estimate has the correct direction;
\item the projective tensor identity with an $\ell_1$-sum is isometric; and
\item the normalized indicator of a finite positive atom spans an
      $L$-summand in $L_1(\mu)$.
\end{itemize}

There is, however, a novelty qualification.  Martín--Perreau--Rueda Zoca
proved the required two-coordinate $\ell_1$-stability for super Daugavet
points, and Lee--Roldán--Tag later restated it explicitly.  Consequently,
Theorem~\ref{thm:supplied} is best regarded as a direct tensor-product
corollary of a known stability theorem.  The supplied proof remains useful
because it is short, self-contained, and makes the tensor mechanism explicit.

\section{Arbitrary $\ell_1$-sums}\label{sec:l1}

The two-coordinate fact extends to arbitrary index sets in exactly the form
needed for tensor applications.

\begin{theorem}[Complete $\ell_1$-sum stability]\label{thm:l1sum}
Let $(E_i)_{i\in I}$ be Banach spaces and
\[
       E=\Bigl(\bigoplus_{i\in I}E_i\Bigr)_1.
\]
Let $z=(z_i)_{i\in I}\in S_E$.  Assume that, whenever $z_i\ne0$, the point
$z_i/\|z_i\|$ is a super Daugavet point of $E_i$.  Then $z$ is a super
Daugavet point of $E$.
\end{theorem}

\begin{proof}
Write
\[
 I_0=\{i\in I:z_i\ne0\},\qquad
 a_i=\|z_i\|,\qquad x_i=z_i/a_i\quad(i\in I_0).
\]
The set $I_0$ is countable and $\sum_{i\in I_0}a_i=1$.

Fix $u=(u_i)_{i\in I}\in B_E$, and put
\[
 s=\sum_{i\notin I_0}\|u_i\|.
\]
Choose numbers $c_i\ge\|u_i\|$ for $i\in I_0$ such that
\[
       \sum_{i\in I_0}c_i=1-s.
\]
This is possible because
$\sum_{i\in I_0}\|u_i\|\le1-s$ and $I_0\ne\varnothing$.

For every $i\in I_0$ with $c_i>0$, use Lemma~\ref{lem:net} to choose a net
$(p_{i,\alpha_i})\subset B_{E_i}$ satisfying
\[
 p_{i,\alpha_i}\xrightarrow{w}u_i/c_i,
 \qquad
 \|x_i-p_{i,\alpha_i}\|\longrightarrow2.
\]
If $c_i=0$, then $u_i=0$ and we simply use the constant value $0$.  On the
product directed set define $w_\alpha=(w_{i,\alpha})_{i\in I}$ by
\[
 w_{i,\alpha}=
 \begin{cases}
   c_i p_{i,\alpha_i},&i\in I_0,\\
   u_i,&i\notin I_0.
 \end{cases}
\]
Then
\[
 \|w_\alpha\|
 \le \sum_{i\in I_0}c_i+
       \sum_{i\notin I_0}\|u_i\|=1,
\]
so $w_\alpha\in B_E$.

We claim that $w_\alpha\to u$ weakly.  Indeed, if
$f=(f_i)\in E^*=(\bigoplus_iE_i^*)_\infty$, convergence holds on every
finite set of coordinates.  Outside a finite set $F\subset I_0$, the tail is
bounded by
\[
  2\|f\|\sum_{i\in I_0\setminus F}c_i,
\]
which can be made arbitrarily small.  This proves weak convergence.

For $i\in I_0$ we have
\begin{align*}
 \|a_ix_i-c_ip_{i,\alpha_i}\|
 &\ge \|x_i-p_{i,\alpha_i}\|-(1-a_i)-(1-c_i),
\end{align*}
so
\[
 \liminf_\alpha\|a_ix_i-c_ip_{i,\alpha_i}\|\ge a_i+c_i.
\]
Therefore, for every finite $F\subset I_0$,
\begin{align*}
 \liminf_\alpha\|z-w_\alpha\|
 &\ge \sum_{i\in F}(a_i+c_i)+s.
\end{align*}
Letting $F$ increase to $I_0$ gives
\[
 \liminf_\alpha\|z-w_\alpha\|
 \ge 1+(1-s)+s=2.
\]
The reverse bound is automatic in the unit ball, hence
$\|z-w_\alpha\|\to2$.  Lemma~\ref{lem:net} proves the assertion.
\end{proof}

\begin{corollary}[Bootstrapping through an $\ell_1$ decomposition]
\label{cor:bootstrap}
Let $Y=(\bigoplus_{i\in I}Y_i)_1$, let $x\in S_X$, and let
$y=(y_i)\in S_Y$.  Suppose that, for every $i$ with $y_i\ne0$,
\[
 x\otimes \frac{y_i}{\|y_i\|}
 \quad\text{is super Daugavet in}\quad X\whotpi Y_i.
\]
Then $x\otimes y$ is super Daugavet in $X\whotpi Y$.
\end{corollary}

\begin{proof}
Use the canonical isometry
\[
 X\whotpi\Bigl(\bigoplus_{i\in I}Y_i\Bigr)_1
 \cong
 \Bigl(\bigoplus_{i\in I}(X\whotpi Y_i)\Bigr)_1
\]
and apply Theorem~\ref{thm:l1sum}.
\end{proof}

\begin{corollary}[The whole scalar $\ell_1$ factor]\label{cor:ell1}
If $x\in S_X$ is a super Daugavet point, then for every set $\Gamma$ and
every $y\in S_{\ell_1(\Gamma)}$, the tensor $x\otimes y$ is a super Daugavet
point of
\[
 X\whotpi\ell_1(\Gamma)\cong\ell_1(\Gamma,X).
\]
\end{corollary}

\begin{proof}
Under the displayed isometry, $x\otimes y$ is the family
$(y_\gamma x)_{\gamma\in\Gamma}$.  Every nonzero normalized coordinate is a
unimodular scalar multiple of $x$, hence is super Daugavet.  Apply
Theorem~\ref{thm:l1sum}.
\end{proof}

This improves the supplied coordinate-vector corollary from $e_\gamma$ to
\emph{every} unit vector of $\ell_1(\Gamma)$.

\section{The whole scalar $L_1$ factor: a later published answer}\label{sec:L1}

The atomic corollary in the supplied packet was substantially superseded by a
later characterization.  Langemets, L\~oo, Mart\'in, Perreau and Rueda Zoca
proved that $f\in S_{L_1(\mu,X)}$ is a super Daugavet point if and only if, on
every atom, either $f$ vanishes or its normalized point value is a super
Daugavet point of $X$ almost everywhere.  Their result immediately gives the
following complete answer for a scalar $L_1$ tensor factor.

\begin{theorem}[Arbitrary scalar $L_1$ vector; published consequence]
\label{thm:L1}
Let $(\Omega,\Sigma,\mu)$ be a measure space, let $X$ be a Banach space, let
$x\in S_X$ be a super Daugavet point, and let $g\in S_{L_1(\mu)}$.  Then
\[
       x\otimes g
\]
is a super Daugavet point of
\[
 X\whotpi L_1(\mu)\cong L_1(\mu,X).
\]
Equivalently, the function $t\mapsto g(t)x$ is a super Daugavet point of
$L_1(\mu,X)$.
\end{theorem}

\begin{proof}
Under the canonical isometry, $x\otimes g$ corresponds to
$f(t)=g(t)x$.  On every atom, either $g=0$ almost everywhere or
\[
   \frac{f(t)}{\|f(t)\|}=\frac{g(t)}{|g(t)|}\,x
\]
almost everywhere.  Unimodular scalar multiples of a super Daugavet point are
again super Daugavet.  The cited characterization of super Daugavet points in
$L_1(\mu,X)$ therefore applies.
\end{proof}

\begin{corollary}[An $L_1$-summand in the second factor]
Suppose $Y=L_1(\mu)\oplus_1 Z$ isometrically.  If $x\in S_X$ is super
Daugavet, then $x\otimes(g,0)$ is super Daugavet in $X\whotpi Y$ for every
$g\in S_{L_1(\mu)}$.
\end{corollary}

\begin{proof}
Combine Theorem~\ref{thm:L1} with the isometry
\[
 X\whotpi Y\cong (X\whotpi L_1(\mu))\oplus_1(X\whotpi Z)
\]
and the two-coordinate $\ell_1$ stability from
Theorem~\ref{thm:supplied}.
\end{proof}

\begin{remark}
Thus the supplied normalized-atom corollary is correct, but it is far from the
largest currently known $L_1$ statement: every unit vector in the scalar
$L_1$ factor works.  This theorem is not claimed as new here; it is a direct
corollary of Proposition~5.8 in the 2026 paper cited above.
\end{remark}

\section{A full positive theorem for ccs Daugavet points}\label{sec:ccs}

A \emph{convex combination of slices} (abbreviated ccs) of $B_X$ is a set
$C=\sum_{i=1}^n\lambda_iS_i$, where each $S_i$ is a slice of $B_X$,
$\lambda_i>0$, and $\sum_i\lambda_i=1$.  A point $x\in S_X$ is a
\emph{ccs Daugavet point} if
\[
       \sup_{u\in C}\|x-u\|=2
\]
for every ccs $C$ of $B_X$.  Bourgain's lemma implies that every ccs Daugavet
point is super Daugavet.  The converse is false in general.

The stronger hypothesis supplies exactly the simultaneous coherence missing
from the general tensor problem.

\begin{theorem}[ccs tensor transfer]\label{thm:ccs}
Let $X$ and $Y$ be Banach spaces over $\K$.  If $x\in S_X$ and $y\in S_Y$
are ccs Daugavet points, then $x\otimes y$ is a ccs Daugavet point of
$X\whotpi Y$.  In particular, $x\otimes y$ is a super Daugavet point.
\end{theorem}

\begin{proof}
Put $E=X\whotpi Y$.  Let
\[
 C=\sum_{i=1}^n\lambda_iS_i
\]
be a ccs of $B_E$, where
\[
 S_i=\{z\in B_E:\Ree B_i(z)>1-\alpha_i\},
 \qquad B_i\in S_{E^*},\quad 0<\alpha_i<1.
\]
We regard $B_i$ as a norm-one bilinear form on $X\times Y$.  Fix $\eta>0$.

For every $i$, choose $a_i\in B_X$ and $b_i\in B_Y$ with
$\Ree B_i(a_i,b_i)>1-\alpha_i/4$, and define the slice
\[
 A_i=\left\{u\in B_X:
 \Ree B_i(u,b_i)>\sup_{r\in B_X}\Ree B_i(r,b_i)-\frac{\alpha_i}{4}\right\}.
\]
Every $u\in A_i$ satisfies $\Ree B_i(u,b_i)>1-\alpha_i/2$.
Since $-x$ is also a ccs Daugavet point, choose $u_i\in A_i$ such that, for
$U=\sum_i\lambda_i u_i$,
\[
       \|x+U\|>2-\eta.
\]
Choose $x^*\in S_{X^*}$, with its phase normalized, so that
$\Ree x^*(x+U)>2-\eta$.  Then
\[
 \Ree x^*(x)>1-\eta,
 \qquad
 A:=\sum_i\lambda_i\bigl(1-\Ree x^*(u_i)\bigr)
   =1-\Ree x^*(U)<\eta.
\]

For every $i$, define the slice
\[
 D_i=\left\{v\in B_Y:
 \Ree B_i(u_i,v)>\sup_{s\in B_Y}\Ree B_i(u_i,s)-\frac{\alpha_i}{2}\right\}.
\]
Every $v\in D_i$ satisfies $\Ree B_i(u_i,v)>1-\alpha_i$.  Use the ccs
Daugavet property of $y$ to choose $v_i\in D_i$ such that, for
$V=\sum_i\lambda_i v_i$,
\[
       \|y-V\|>2-\eta.
\]
Choose $y^*\in S_{Y^*}$, again with normalized phase, so that
$\Ree y^*(y-V)>2-\eta$.  Then
\[
 \Ree y^*(y)>1-\eta,
 \qquad
 D:=\sum_i\lambda_i\bigl(1+\Ree y^*(v_i)\bigr)
   =1+\Ree y^*(V)<\eta.
\]

Set $w=\sum_i\lambda_i u_i\otimes v_i$.  Since
$u_i\otimes v_i\in S_i$, we have $w\in C$.  Put
$p_i=x^*(u_i)$ and $q_i=y^*(v_i)$.  Since $|p_i|,|q_i|\le1$,
\[
 \sum_i\lambda_i|p_i-1|^2\le2A,
 \qquad
 \sum_i\lambda_i|q_i+1|^2\le2D.
\]
Cauchy--Schwarz therefore gives
\begin{align*}
 \left|1+\sum_i\lambda_i p_iq_i\right|
 &\le \sum_i\lambda_i|p_i-1|+
       \sum_i\lambda_i|q_i+1|\\
 &\le \sqrt{2A}+\sqrt{2D}<2\sqrt{2\eta}.
\end{align*}
Also, if $p=x^*(x)$ and $q=y^*(y)$, then
$|p-1|,|q-1|<\sqrt{2\eta}$, and hence
$\Ree(pq)>1-2\sqrt{2\eta}$.

The norm-one bilinear form $\Phi(r,s)=x^*(r)y^*(s)$ now yields
\begin{align*}
 \|x\otimes y-w\|
 &\ge \Ree\Phi(x\otimes y-w)\\
 &>2-4\sqrt{2\eta}.
\end{align*}
Letting $\eta\downarrow0$ proves that $x\otimes y$ is ccs Daugavet.
Bourgain's lemma gives the super Daugavet conclusion.
\end{proof}

\begin{corollary}[Two $L_1$-preduals]
Let $X$ and $Y$ be $L_1$-preduals.  If $x\in S_X$ and $y\in S_Y$ are
super Daugavet points, then $x\otimes y$ is a super Daugavet point of
$X\whotpi Y$.
\end{corollary}

\begin{proof}
In an $L_1$-predual, the standard diametral notions under discussion,
including super Daugavet and ccs Daugavet, coincide.  Apply
Theorem~\ref{thm:ccs}.
\end{proof}

\begin{corollary}
If Banach spaces $X$ and $Y$ both have the Daugavet property, then every
elementary tensor $x\otimes y$ with $x\in S_X$ and $y\in S_Y$ is a ccs
Daugavet point of $X\whotpi Y$.
\end{corollary}

\begin{remark}
The last corollary concerns elementary tensors.  It does not assert that the
whole projective tensor product has the Daugavet property; that broader
space-level tensor question remains delicate and is known under additional
hypotheses such as the weak operator Daugavet property.
\end{remark}

\section{The elementary-target local theorem}\label{sec:elementary}

Even without ccs hypotheses, the desired conclusion can be proved in every
weak neighborhood whose center is an elementary tensor.

\begin{theorem}[Local transfer at elementary targets]
\label{thm:elementary-local}
Let $X$ and $Y$ be Banach spaces over $\K$, and let $x\in S_X$ and
$y\in S_Y$ be super Daugavet points.  Put $p=x\otimes y$.  Let $a\in B_X$,
$b\in B_Y$, and let $W$ be a relatively weak neighborhood of $a\otimes b$
in $B_{X\whotpi Y}$.  Then, for every $\varepsilon>0$, there are $u\in B_X$
and $v\in B_Y$ such that
\[
       u\otimes v\in W
       \qquad\text{and}\qquad
       \|p-u\otimes v\|>2-\varepsilon.
\]
\end{theorem}

\begin{proof}
Choose a basic relatively weak neighborhood $W_0\subset W$ of $a\otimes b$
of the form
\[
 W_0=\left\{z\in B_{X\whotpi Y}:
 |B_j(z-a\otimes b)|<\delta,
 \ 1\le j\le m\right\},
\]
where the $B_j$ are bounded bilinear forms.  After rescaling the finitely many
forms, we may assume $\|B_j\|\le1$.  Choose $\eta>0$ so small that
$2\eta<\delta$ and $4\sqrt{2\eta}<\varepsilon$.

The set
\[
 U=\{u\in B_X:|B_j(u-a,b)|<\eta\text{ for }1\le j\le m\}
\]
is a nonempty relatively weakly open subset of $B_X$.  Since $-x$ is super
Daugavet, choose $u\in U$ with $\|x+u\|>2-\eta$.  Now
\[
 V=\{v\in B_Y:|B_j(u,v-b)|<\eta\text{ for }1\le j\le m\}
\]
is a nonempty relatively weakly open subset of $B_Y$.  Since $y$ is super
Daugavet, choose $v\in V$ with $\|y-v\|>2-\eta$.
For every $j$,
\begin{align*}
 |B_j(u,v)-B_j(a,b)|
 &\le |B_j(u,v-b)|+|B_j(u-a,b)|\\
 &<2\eta<\delta,
\end{align*}
so $u\otimes v\in W_0\subset W$.

Choose $x^*\in S_{X^*}$ and $y^*\in S_{Y^*}$, with phases normalized, such
that
\[
 \Ree x^*(x+u)>2-\eta,
 \qquad
 \Ree y^*(y-v)>2-\eta.
\]
Thus each of $x^*(x)$, $x^*(u)$ and $y^*(y)$ is within
$\sqrt{2\eta}$ of $1$, while $y^*(v)$ is within $\sqrt{2\eta}$ of $-1$.
For the norm-one bilinear form $\Phi(r,s)=x^*(r)y^*(s)$, it follows that
\[
 \Ree\Phi(x\otimes y-u\otimes v)>2-4\sqrt{2\eta}>2-\varepsilon.
\]
Hence $\|x\otimes y-u\otimes v\|>2-\varepsilon$.
\end{proof}

\begin{corollary}
If every nonempty relatively weakly open subset of $B_{X\whotpi Y}$ contains
an elementary tensor, then super Daugavet points $x\in S_X$ and $y\in S_Y$
produce a super Daugavet tensor $x\otimes y$.
\end{corollary}

The hypothesis of this corollary is generally false in the strongest possible
way.  Petitjean proved that the set
\[
       \{a\otimes b:a\in X,\ b\in Y\}
\]
is weakly closed in $X\whotpi Y$.  Thus every non-elementary tensor has a
weak neighborhood containing no elementary tensor.  This shows precisely why
an argument based only on selecting one elementary tensor from the target weak
open set cannot settle the general problem.

\section{Why the remaining step is genuinely different}\label{sec:obstruction}

The known and new information can be summarized as follows.

\begin{center}
\begin{tabular}{p{0.31\textwidth}p{0.57\textwidth}}
\toprule
Input & Known conclusion for $x\otimes y$ \\
\midrule
$x$ super $\Delta$, $y$ arbitrary & super $\Delta$ \\
$x,y$ Daugavet & Daugavet \\
$x,y$ ccs Daugavet & ccs Daugavet, hence super Daugavet \\
$x$ super Daugavet, $y\in L_1(\mu)$ arbitrary & super Daugavet \\
$x,y$ super Daugavet, target weak neighborhood contains an elementary tensor
& an elementary tensor in that neighborhood can be chosen almost antipodal to
$x\otimes y$ \\
$x,y$ super Daugavet, arbitrary target weak neighborhood & open in general \\
\bottomrule
\end{tabular}
\end{center}

Two tempting shortcuts fail.

\paragraph{The two known transfers do not combine automatically.}
From existing results, if $x$ and $y$ are super Daugavet, then $x\otimes y$ is
both a Daugavet point and a super $\Delta$-point.  There is no general theorem
identifying the conjunction of those two pointwise properties with the super
Daugavet property: they control slices and weak neighborhoods containing the
point, respectively, rather than all relatively weakly open subsets.  Thus the
two conclusions alone do not settle the question.

\paragraph{Finite convex decompositions require coherence.}
Every element of $B_{X\whotpi Y}$ is weakly approximable by convex
combinations of elementary tensors.  Perturbing each factor separately using
the super Daugavet property gives individual norming functionals, but these
functionals need not be compatible across the different summands.  The proof
of Theorem~\ref{thm:ccs} succeeds because the ccs hypotheses create common
functionals for the two weighted families.  A super Daugavet point need not be
a ccs $\Delta$-point, so this coherence cannot be assumed in the general
problem.

This obstruction is consistent with recent work on weakly open sets in tensor
products: slices always contain elementary tensors, while arbitrary weakly
open subsets need not.  That distinction is identified in the recent
literature as the central difficulty in transferring diameter-two and
Daugavet-type properties to projective tensor products.

\section{Literature search and status as of 21 June 2026}

The search included the exact sentence from the source paper, combinations of
``super Daugavet'', ``projective tensor product'', and ``elementary tensor'',
the papers citing the source that were identifiable in bibliographic and
full-text searches, the 2024 monograph on diametral notions, the 2025 article
on vector-valued function spaces, the 2026 paper on the super alternative
Daugavet property, a 2025 doctoral thesis on Daugavet and $\Delta$ points, the
2025 monograph manuscript on the Daugavet property, and 2024--2026 papers on
weakly open sets and elementary tensors in projective tensor products.

No paper located in that search states a proof or a counterexample to
\eqref{eq:main-question}.  The most relevant later developments instead give:
\begin{itemize}[leftmargin=1.8em]
\item a complete characterization of super Daugavet points in
      $L_1(\mu,X)$, which settles every scalar $L_1$ tensor factor;
\item positive space-level tensor results under the weak operator Daugavet
      property;
\item a precise explanation of the weak-open-set obstruction; and
\item weak closedness of the elementary tensors.
\end{itemize}
Accordingly, the responsible status statement is: \emph{the general question
appears to remain open in the literature found through 21 June 2026}.  The
ccs-transfer theorem and the elementary-target theorem in this report were not
located in that literature; they should nevertheless be regarded as new
proofs proposed here, not as peer-reviewed results.

\section{Conclusion}

The supplied partial result is mathematically correct, but its principal
stability lemma is already known.  The strongest clean package established or
identified here is:
\begin{enumerate}[label=(\arabic*),leftmargin=2.2em]
\item arbitrary $\ell_1$-sums preserve a point assembled from super Daugavet
      coordinates;
\item the 2026 $L_1(\mu,X)$ characterization implies that $x\otimes g$ is
      super Daugavet for every $g\in S_{L_1(\mu)}$ whenever $x$ is;
\item two ccs Daugavet points have a ccs Daugavet elementary tensor, over both
      real and complex scalars; and
\item two super Daugavet points solve the tensor problem in every weak
      neighborhood centered at an elementary tensor.
\end{enumerate}
The unresolved part is exactly the passage to weak neighborhoods around
non-elementary tensors, where no rank-one representative is available and the
separate factor perturbations need not admit a common norming bilinear form.

\begin{thebibliography}{99}

\bibitem{HLPV}
R.~Haller, J.~Langemets, Y.~Perreau, and T.~Veeorg,
\emph{Unconditional bases and Daugavet renormings},
J. Funct. Anal. \textbf{286} (2024), no.~12, Paper No.~110421, 31 pp.,
\href{https://doi.org/10.1016/j.jfa.2024.110421}{doi:10.1016/j.jfa.2024.110421}.

\bibitem{DantasJungRZ}
S.~Dantas, M.~Jung, and A.~Rueda Zoca,
\emph{Daugavet points in projective tensor products},
Q. J. Math. \textbf{73} (2022), 443--459.

\bibitem{MPRZ}
M.~Martín, Y.~Perreau, and A.~Rueda Zoca,
\emph{Diametral notions for elements of the unit ball of a Banach space},
Dissertationes Math. \textbf{594} (2024), 1--61,
\href{https://doi.org/10.4064/dm230728-21-3}{doi:10.4064/dm230728-21-3}.

\bibitem{LRT}
H.~J. Lee, Ó.~Roldán, and H.-J. Tag,
\emph{On various diametral notions of points in the unit ball of some
vector-valued function spaces},
Rev. R. Acad. Cienc. Exactas Fís. Nat. Ser. A Mat. RACSAM
\textbf{119} (2025), no.~4, Paper No.~88,
\href{https://doi.org/10.1007/s13398-025-01757-6}{doi:10.1007/s13398-025-01757-6}.

\bibitem{SuperADP}
J.~Langemets, M.~L\~oo, M.~Mart\'in, Y.~Perreau, and A.~Rueda Zoca,
\emph{The super Alternative Daugavet property for Banach spaces},
Proc. Edinb. Math. Soc. \textbf{69} (2026), 471--499,
\href{https://doi.org/10.1017/S0013091525101168}{doi:10.1017/S0013091525101168}.

\bibitem{RZweakopen}
A.~Rueda Zoca,
\emph{Weak operator Daugavet property and weakly open sets in tensor product
spaces},
Rev. Mat. Complut. \textbf{39} (2026), 149--164,
\href{https://doi.org/10.1007/s13163-025-00531-x}{doi:10.1007/s13163-025-00531-x}.

\bibitem{Petitjean}
C.~Petitjean,
\emph{The set of elementary tensors is weakly closed in projective tensor
products},
Bull. Aust. Math. Soc. \textbf{111} (2025), 336--343,
\href{https://doi.org/10.1017/S0004972724000376}{doi:10.1017/S0004972724000376}.

\bibitem{Shvidkoy}
R.~V. Shvidkoy,
\emph{Geometric aspects of the Daugavet property},
J. Funct. Anal. \textbf{176} (2000), 198--212.

\bibitem{Ryan}
R.~A. Ryan,
\emph{Introduction to Tensor Products of Banach Spaces},
Springer Monographs in Mathematics, Springer, London, 2002.

\bibitem{KadetsBook}
V.~Kadets, M.~Martín, A.~Rueda Zoca, and D.~Werner,
\emph{Banach Spaces with the Daugavet Property}, book manuscript, 2025.

\end{thebibliography}

\end{document}
